%%%%%%%%%%%%%%%%%%%%%%%%%%%%%%%%%%%%%%% Titelseite

% LaTex will Bilder als EPS haben, damit es richtig kompiliert wird. Deswegen "imagemagick" herunter laden und mit "mogrify -format eps *.jpg" alle Bilder umcodieren, sodass sie auch im eps-Format vorliegen.

\thispagestyle{empty}
$\begin{minipage}{0.65\textwidth}
\hspace{8mm}
\includegraphics[scale=0.21]{Bilder/1} \\
Otto-von-Guericke-Universität \\
\makebox[\widthof{Otto-von-Guericke-Universität}] [c] {Magdeburg}
\end{minipage}
%\hspace{18mm}
\begin{minipage}{0.3\textwidth}
\vspace{2mm}
\includegraphics[scale=0.27]{Bilder/2} \\
\begin{center}
\vspace{-2mm}
Eingebettete Systeme\\
\makebox [\widthof{Eingebettete Systeme}][c]{und Betriebssysteme}
\end{center}
\end{minipage}$\\[2cm]

\begin{center}
{\bf {\huge Masterarbeit}} \\[1cm] 
{\LARGE{Lighthouse Keeper - Robotergestützte Erfassung und Evaluation von iBeacon Konfigurationen}} \par\medskip
von \par\smallskip
{\bf André Pieper \\
Geb. 28.03.1988 in Berlin \\
Matrikelnummer: 184960}\\ [1cm]
\end{center}

\begin{tabular}{ll}
Erstprüfer: & Junior-Prof. Dr. Ing. Sebastian Zug \\
& Otto-von-Guericke-Universität \\
& Fakultät für Informatik \\
& Institut für Verteilte Systeme \\
& Universitätsplatz 2, D-39106, Magdeburg \\ \\

Zweitprüfer: & Prof. Dr. Abbas Omar \\
& Otto-von-Guericke-Universität \\
& Fakultät für Elektrotechnik und Informationstechnik \\
& Institut für Informations- und Kommunikationstechnik\\
& Universitätsplatz 2, D-39106, Magdeburg \\ \\
Betreuer: & Junior-Prof. Dr. Ing. Sebastian Zug 
\end{tabular}

\newpage%%%%%%%%%%%%%%%%%%%%% Seite 2

\thispagestyle{empty}
\section*{}

\newpage%%%%%%%%%%%%%%%%%%%%% Seite 4

\pagenumbering{Roman}
\tableofcontents

\newpage%%%%%%%%%%%%%%%%%%%%% Seite 5

\chapter*{Vorwort}
\addcontentsline{toc}{chapter}{Vorwort}
Hier steht dann das Vorwort. Hier stehen so Sachen wie, dass iBeacons die große Revolution in der Indoor-Lokalisierung sind. Dass sie in Flughäfen, Apple-Geschäften und (weitere Einsatzgebiete)... schon eingesetzt werden. Jedoch fehlt es an einer geeigneten Strategie die iBeacons zu verteilen (bisher wurde es abgeschätzt) und die Planung von Aufbau und Instandsetzumg von iBeacon-Lokalisierungsfeldern optimal umzusetzen. Denn in großen Gebieten, wie z.B. Einkaufshäuser, Lagerhallen und Flughäfen braucht man Hunderte von iBeacons, um die gesamte Fläche zu 100 Prozent abzudecken. Im Idealfall kann ich später in Matlab noch Gebiete einpflegen, an der die Karte nach Regionen aufgeteilt wird, in denen die Genauigkeit der Lokalisierung eingestellt werden kann. Damit ließe sich Robotergestützt ein Gebiet abfahren, in denen Vorschläge für die Positionierung von iBeacons anhand eines Optimierungsalgorithmusses generiert werden. Hierbei ließe sich wunderschön eine Pareto-Front bilden, die aus Anzahl von iBeacons und Genauigkeit pro Region Positionierungsvorschläge entwirft. Je nach Kundenwunsch.  

\newpage%%%%%%%%%%%%%%%%%%%%% Seite 5

\listoffigures
\addcontentsline{toc}{chapter}{Abbildungsverzeichnis}

\newpage%%%%%%%%%%%%%%%%%%%%% Seite 6

\chapter*{Abkürzungsverzeichnis}
\addcontentsline{toc}{chapter}{Abkürzungsverzeichnis}

KF - Kalmanfilter\\


\newpage%%%%%%%%%%%%%%%%%%%%% Seite 8